\documentclass[13pt, a4paper]{article}

\usepackage{amsmath}
\usepackage{amssymb}
\usepackage{amsthm}
\usepackage{bussproofs}
\usepackage{fancyhdr}
\usepackage[margin=3cm]{geometry}
\usepackage{mathpartir}
\usepackage{stmaryrd}

\newtheorem{theorem}{Theorem}
\newtheorem{lemma}{Lemma}
\newtheorem{proposition}{Proposition}

\pagestyle{fancy}
\fancyhf{} % Clear all header and footer fields

% Header: left, center, right
\fancyhead[L]{Programming Languages 2}
\fancyhead[C]{Homework 04}
\fancyhead[R]{\today}
\fancyfoot[C]{\thepage}

\begin{document}

Simon Sutoris 7542170

%===============================================================================
\section*{1 Typing}
%===============================================================================

\subsection*{1.} %==============================================================

\begin{prooftree}
  \AxiomC{}
  \UnaryInfC{\texttt{Fl}(1) : \texttt{Fl}}
  \AxiomC{$\lightning$}
  \UnaryInfC{\texttt{pow}(\texttt{Fl}(5), \texttt{Fl}(0.5)) : \texttt{Fl}}
  \BinaryInfC{\texttt{Fl}(1) + \texttt{pow}(\texttt{Fl}(5), \texttt{Fl}(0.5)) : \texttt{Fl}}
  \AxiomC{}
  \UnaryInfC{\texttt{Fl}(2) : \texttt{Fl}}
  \BinaryInfC{(\texttt{Fl}(1) + \texttt{pow}(\texttt{Fl}(5), \texttt{Fl}(0.5))) / \texttt{Fl}(2) : \texttt{Fl}}
\end{prooftree}

The term is ill-typed, as \textsc{Ty-Pow} requires $t_2 : \texttt{Int}$, but $t_2 = \texttt{Fl}(0.5) : \texttt{Fl}$.

\subsection*{2.} %==============================================================

\begin{prooftree}
  \AxiomC{}
  \UnaryInfC{\texttt{Fl}(3) : \texttt{Fl}}
  \AxiomC{}
  \UnaryInfC{\texttt{Int}(4) : \texttt{Int}}
  \UnaryInfC{\texttt{AsFloat}(\texttt{Int}(4)) : \texttt{Fl}}
  \BinaryInfC{\texttt{Fl}(3)\ast\texttt{AsFloat}(\texttt{Int}(4)) : \texttt{Fl}}
  \UnaryInfC{\texttt{round}(\texttt{Fl}(3)\ast\texttt{AsFloat}(\texttt{Int}(4))) : \texttt{Int}}
\end{prooftree}

\subsection*{3.} %==============================================================

\begin{prooftree}
  \AxiomC{}
  \UnaryInfC{\texttt{Int}(2) : \texttt{Int}}
  \UnaryInfC{texttt{AsFloat}(\texttt{Int}(2)) : \texttt{Fl}}
  \AxiomC{}
  \UnaryInfC{\texttt{Int}(64) : Int}
  \BinaryInfC{\texttt{pow}(\texttt{AsFloat}(\texttt{Int}(2)), \texttt{Int}(64)) : \texttt{Fl}}
\end{prooftree}

\subsection*{4.} %==============================================================

\begin{prooftree}
  \AxiomC{$\lightning$}
  \UnaryInfC{\texttt{Fl}(4)/\texttt{Int}(2) : $\lightning$}
\end{prooftree}

The term is ill-typed, as \textsc{Ty-BinOp} requires the types $t_1 : \tau$ and $t_2 : \tau$ to be equal, but $t_1 = \texttt{Fl}(4) : \texttt{Fl}$ has a different type than $t_2 = \texttt{Int}(2) : \texttt{Int}$

%===============================================================================
\section*{2 Progress and Preservation}
%===============================================================================

\subsection*{Theorem 1 \textmd{Preservation}} %=================================

\subsubsection*{\textmd{\textsc{Ty-Int}}}
Term is a value. $\square$

\subsubsection*{\textmd{\textsc{Ty-BinOp}, \textsc{Ty-Float}, \textsc{Ty-Round}, \textsc{Ty-ToFloat}, \textsc{Ty-Pow}}}
Analog to \texttt{Ty-Int}.

\subsubsection*{\textmd{\textsc{E-FlOp}}}
\begin{align}
t_1 &= \texttt{Fl}(f_1) \oplus \texttt{Fl}(f_2) : \texttt{Fl} \\
t_2 &= \texttt{Fl}(f_1 \oplus f_2) : \texttt{Fl}
\end{align}
Types of $t_1$ and $t_2$ match. $\square$

\subsubsection*{\textmd{\textsc{E-IntOp}}}
\begin{align}
t_1 &= \texttt{Int}(i_1) \oplus \texttt{Int}(i_2) : \texttt{Int} \\
t_2 &= \texttt{Int}(i_1 \oplus i_2) : \texttt{Int}
\end{align}
Types of $t_1$ and $t_2$ match. $\square$

\subsubsection*{\textmd{\textsc{E-RoundFl}}}
\begin{align}
t_1 &= \texttt{round}(\texttt{Fl}(f)) : \texttt{Int} \\
t_2 &= \texttt{Int}(\lfloor f\rceil) : \texttt{Int}
\end{align}
Types of $t_1$ and $t_2$ match. $\square$

\subsubsection*{\textmd{\textsc{E-IntFl}}}
\begin{align}
t_1 &= \texttt{AsFloat}(\texttt{Int}(i)) : \texttt{Fl} \\
t_2 &= \texttt{Fl}(i) : \texttt{Fl}
\end{align}
Types of $t_1$ and $t_2$ match. $\square$

\subsubsection*{\textmd{\textsc{E-PowFl}}}
\begin{align}
t_1 &= \texttt{pow}(\texttt{Fl}(f),\texttt{Int}(i)) : \texttt{Fl} \\
t_2 &= \texttt{Fl}(f^i) : \texttt{Fl}
\end{align}
Types of $t_1$ and $t_2$ match. $\square$

\subsubsection*{\textmd{\textsc{E-PowInt}}}
\begin{align}
t_1 &= \texttt{pow}(\texttt{Int}(i_1),\texttt{Int}(i_2)) : \texttt{Fl} \\
t_2 &= \texttt{Fl}(i_1^{i_2}) : \texttt{Fl}
\end{align}
Types of $t_1$ and $t_2$ match. $\square$

\subsubsection*{\textmd{\textsc{E-BinOp1}}}
Due to \textsc{Ty-BinOp}, $\oplus$ has the same type as its operands. Due to \textsc{E-FlOp}, \textsc{E-IntOp}, \textsc{E-RoundFl}, \textsc{E-IntFL}, \textsc{E-PowFl} and \textsc{E-PowInt}, $t_1$ and $t_1'$ in $t_1 \mapsto t_2$ have the same type $\tau$. Therefore $t_1 \oplus t_2$ and $t_1' \oplus t_2$ have the same type $\tau$. $\square$

\subsubsection*{\textmd{\textsc{E-BinOp2}, \textsc{E-Round}, \textsc{E-AsFloat}, \textsc{E-Pow1}, \textsc{E-Pow2}}}
Analog to \textsc{E-BinOp1}.

\noindent$\blacksquare$

\subsection*{Theorem 2 \textmd{Progress}} %=====================================

\subsubsection*{\textmd{\textsc{Ty-Int}, \textsc{Ty-BinOp}, \textsc{Ty-Float}, \textsc{Ty-Round}, \textsc{Ty-ToFloat}, \textsc{Ty-Pow}}}

These rules are all values. $\square$

\subsubsection*{\textmd{\textsc{E-FlOp}, \textsc{E-IntOp}, \textsc{E-RoundFl}, \textsc{E-IntFl}, \textsc{E-PowFl}, \textsc{E-PowInt}}}

These rules are all reducible terms. $\square$

\subsubsection*{\textmd{\textsc{E-BinOp1}}}

\textbf{IA}: $t_1$ is a value, i.e. not reducible: $t_1 =t_1'$. Therefore $t_1 \oplus t_2 = t_1' \oplus t_2$.

\noindent \textbf{IS}: $t_1$ is reducible. $\square$

\subsubsection*{\textmd{\textsc{E-BinOp2}, \textsc{E-Round}, \textsc{E-AsFloat}, \textsc{E-Pow1}, \textsc{E-Pow2}}}
Analog to \textsc{E-BinOp1}.

\noindent$\blacksquare$

%===============================================================================
\section*{3 Designing a type system}
%===============================================================================

\subsection*{1.} %==============================================================

\begin{center}
\[
\begin{array}{lrr}
t & := & \\
& & \texttt{Int}(n) \\
& | & \texttt{List}(\texttt{Nil}) \\
& | & \texttt{Cons}(t, t) \\
& | & \texttt{tail}(t, t) \\
& | & \texttt{succ}(t, t) \\
& | & \texttt{pred}(t, t) \\
\end{array}
\]
\end{center}

\begin{prooftree}
  \AxiomC{$x:T \in \Gamma$}
  \RightLabel{\textsc{Ty-Var}}
  \UnaryInfC{\texttt{Int}(n): \texttt{Int}}
\end{prooftree}

\begin{prooftree}
  \AxiomC{}
  \RightLabel{\textsc{Ty-Nil}}
  \UnaryInfC{\texttt{List}(\texttt{Nil}): \texttt{List}}
\end{prooftree}

\begin{prooftree}
  \AxiomC{$t_1$ : \texttt{Int}}
  \AxiomC{$t_2$ : \texttt{List}}
  \RightLabel{\textsc{Ty-Cons}}
  \BinaryInfC{\texttt{Cons}($t_1$, $t_2$): \texttt{List}}
\end{prooftree}

\begin{prooftree}
  \AxiomC{t : \texttt{List}}
  \RightLabel{\textsc{Ty-Tail}}
  \UnaryInfC{\texttt{tail}(t): \texttt{List}}
\end{prooftree}

\begin{prooftree}
  \AxiomC{t : \texttt{Int}}
  \RightLabel{\textsc{Ty-Succ}}
  \UnaryInfC{\texttt{succ}(t): \texttt{Int}}
\end{prooftree}

\begin{prooftree}
  \AxiomC{t : \texttt{Int}}
  \RightLabel{\textsc{Ty-Pred}}
  \UnaryInfC{\texttt{pred}(t): \texttt{Int}}
\end{prooftree}

\subsection*{2.} %==============================================================

\begin{center}
\[
\begin{array}{rcll}
\texttt{tail}(\texttt{Cons}(\texttt{Int}(n), \texttt{List}(vl))) & \mapsto & \texttt{List}(vl) & \textsc{Ty-TailCons} \\
\texttt{tail}(\texttt{List}(\texttt{Nil})) & \mapsto & \texttt{List}(\texttt{Nil}) & \textsc{Ty-TailNil} \\
\texttt{succ}(\texttt{Int}(n)) & \mapsto & \texttt{Int}(n + 1) & \textsc{Ty-SuccOp} \\
\texttt{succ}(\texttt{Pred}(n)) & \mapsto & \texttt{Int}(n - 1) & \textsc{Ty-PredOp} \\
\end{array}
\]
\end{center}

\subsection*{3. Preservation} %=================================================

\subsubsection*{\textmd{\textsc{Ty-Int}}}
Term is a value. $\square$

\subsubsection*{\textmd{\textsc{Ty-Nil}, \textsc{Ty-Cond}, \textsc{Ty-Tail}, \textsc{Ty-Succ}, \textsc{Ty-Pred}}}
Analog to \texttt{Ty-Int}.

\subsubsection*{\textmd{\textsc{Ty-TailCons}}}
\begin{align}
t_1 &= \texttt{Cons}(\texttt{Int}(n), \texttt{List}(vl))) : \texttt{List} \\
t_2 &= \texttt{List}(vl) : \texttt{List}
\end{align}
Types of $t_1$ and $t_2$ match. $\square$

\subsubsection*{\textmd{\textsc{Ty-TailNil}}}
\begin{align}
t_1 &= \texttt{tail}(\texttt{List}(\texttt{Nil})) : \texttt{List} \\
t_2 &= \texttt{List}(\texttt{Nil}) : \texttt{List}
\end{align}
Types of $t_1$ and $t_2$ match. $\square$

\subsubsection*{\textmd{\textsc{Ty-SuccOp}}}
\begin{align}
t_1 &= \texttt{succ}(\texttt{Int}(n)) : \texttt{Int} \\
t_2 &= \texttt{Int}(n + 1) : \texttt{List}
\end{align}
Types of $t_1$ and $t_2$ match. $\square$

\subsubsection*{\textmd{\textsc{Ty-PredOp}}}
\begin{align}
t_1 &= \texttt{pred}(\texttt{Int}(n)) : \texttt{Int} \\
t_2 &= \texttt{Int}(n - 1) : \texttt{List}
\end{align}
Types of $t_1$ and $t_2$ match. $\square$

\subsubsection*{\textmd{\textsc{E-Cons1}}}
Due to \textsc{Ty-TailCons}, \textsc{Ty-TailNil}, \textsc{Ty-SuccOp} and \textsc{Ty-PredOp}, $t_1$ and $t_1'$ in $t_1 \mapsto t_2$ have the same type $\tau$. Therefore $\texttt{Cons}(t_1, t_2)$ and $\texttt{Cons}(t_1', t_2)$ have the same type $\tau$. $\square$

\subsubsection*{\textmd{\textsc{E-Cons2}, \textsc{E-Tail}, \textsc{E-Succ}, \textsc{E-Pred}}}
Analog to \textsc{E-Cons1}.

\noindent$\blacksquare$

\subsection*{3. Progress} %=====================================================

\subsubsection*{\textmd{\textsc{Ty-Int}, \textsc{Ty-Nil}, \textsc{Ty-Cons}, \textsc{Ty-Tail}, \textsc{Ty-Succ}, \textsc{Ty-Pred}}}

These rules are all values. $\square$

\subsubsection*{\textmd{\textsc{Ty-TailCons}, \textsc{Ty-TailNil}, \textsc{Ty-SuccOp}, \textsc{Ty-PredOp}}}

These rules are all reducible terms. $\square$

\subsubsection*{\textmd{\textsc{E-Cons1}}}

\textbf{IA}: $t_1$ is a value, i.e. not reducible: $t_1 =t_1'$. Therefore $\texttt{Cons}(t_1, t_2) = \texttt{Cons}(t_1', t_2)$.

\noindent \textbf{IS}: $t_1$ is reducible. $\square$

\subsubsection*{\textmd{\textsc{E-Cons2}, \textsc{E-Tail}, \textsc{E-Succ}, \textsc{E-Pred}, \textsc{E-Pow2}}}
Analog to \textsc{E-Cons1}.

\noindent$\blacksquare$

%===============================================================================
\section*{4 Types and Programming Languages}
%===============================================================================

No questions.

%===============================================================================
\section*{5 Nameless Representations}
%===============================================================================

\begin{enumerate}
\item $\lambda.\lambda.1\,0$
\item $\lambda.2\,(\lambda.\lambda.1)$
\item $\lambda.((\lambda.2)\,(\lambda.0))$
\end{enumerate}

\end{document}
